\begin{abstract}
Self-improving agent harnesses such as the Darwin G\"odel Machine (DGM) let a coding agent edit its own scaffolding, and
they keep the edits that score better and report large gains on coding benchmarks. In three pre-registered rounds, we
compare what such loops find with standard engineering changes (``textbook moves''), one-shot redesigns by a language
model, and placebo loops that make no edits; we also re-run released agents of the DGM and of ADAS with newer
models. On hard LiveCodeBench problems with Claude Haiku~5.5, the starting agent often stops on an empty turn because its
output-token limit is too small for the model's reasoning. A single textbook move that raises this limit recovers about 90\% of the
gain of the best loop agent and costs less than half as much per task. Four of five one-shot redesigns by Claude
Opus~5.5 find the same fix without any evaluation. The best loop agent still leads the textbook moves by 1--5 percentage
points in pass rate. On held-out SWE-bench Verified tasks, run in the DGM's harness, the DGM's released best agent does not outperform
its initial agent; three textbook moves on the initial agent outperform the best agent at a quarter of its cost, whether
or not the patches of runs that ended early, most of them the best agent's, are kept. Claims of self-improvement should be measured against
textbook moves, one-shot redesigns, placebo loops, held-out tasks and newer models.
\end{abstract}
